% Lengkapi identitas kelompok sebelum pengumpulan.
\newcommand{\JenisLaporan}{LAPORAN TUGAS 2}
\newcommand{\KodeMK}{IF4020}
\newcommand{\NamaMK}{Kriptografi}
\newcommand{\JudulLaporan}{Perancangan dan Analisis Brainrot Cipher}
\newcommand{\NamaAlgoritma}{Brainrot Cipher}
\newcommand{\NomorKelompok}{11}
\newcommand{\NamaSatu}{Muhammad Ra'if Alkautsar}
\newcommand{\NIMSatu}{13523011}
\newcommand{\NamaDua}{Muhammad Aditya Rahmadeni}
\newcommand{\NIMDua}{13523028}
\newcommand{\NamaTiga}{Fityatul Haq}
\newcommand{\NIMTiga}{13523116}
\newcommand{\NamaEmpat}{[Nama anggota 4]}
\newcommand{\NIMEmpat}{[NIM 4]}
\newif\ifEmpatAnggota
\EmpatAnggotafalse % Ganti dengan \EmpatAnggotatrue untuk kelompok 4 orang.
\newcommand{\FotoKelompok}{foto-kelompok.png}
\newcommand{\ProgramStudi}{PROGRAM STUDI TEKNIK INFORMATIKA}
\newcommand{\Fakultas}{SEKOLAH TEKNIK ELEKTRO DAN INFORMATIKA}
\newcommand{\Institusi}{INSTITUT TEKNOLOGI BANDUNG}
\newcommand{\AlamatInstitusi}{JL. GANESHA 10, BANDUNG 40132}
\newcommand{\TahunAkademik}{2026/2027}
\newcommand{\Repositori}{https://github.com/mraifalkautsar/Tugas2-IF4020-2026-K11}
% Tautan bonus diisi langsung di bagian/8_lampiran.tex jika tersedia.
